% sup3_00_intro: Supplement part 3, section opening and conventions.
% Written proofs of the cube steps of Ba3 Prop. 1, p. 432-433: the component cube (S11), the properties of the cube
% split that the later steps use (S12, printed), the constant field B_0 (S13), the expansion in the new background and
% its classes (S14), and the localisation of the old-vertex components (S15). Input from SUPPLEMENT.tex through
% INPUTS_part3.tex.

\section{The sketched steps of Proposition~1, III: the component cube, the constant field, the expansion in the new background and the old-vertex components}\label{sup3:sec}

This part contains the steps of \cite[pp.~432--433]{Ba3} that prepare the passage from the domain $\Omega$ to a cube $\square$ of constant side and from the background $\tilde B$ to a constant background $B_0$:

\begin{center}
\begin{tabular}{lll}
step & \cite{Ba3} page & printed object\\\hline
S11 & 432 & ``the localizations of each connected component are contained in some cube of a linear size depending on\\
 & & a number of vertices, hence on $\bar n$ only''\\
S12 & 433 & the cube split $G_k(\Omega,\tilde B)=G_k(\square,\tilde B)+\delta G_k(\Omega,\square,\tilde B)$ (printed; only its three properties (P-i)--(P-iii) are used)\\
S13 & 433 & the constant field $B_0$ with $|\tilde B-B_0|\le O(r(L^k\eps)p(L^k\eps))$ on $\square$\\
S14 & 433 & the expansion in $\tilde B'=\tilde B-B_0$ by \cite[(3.14), (3.44), (3.45)]{Ba1}, and its decomposition into renormalised classes\\
S15 & 433 & ``If a class contains an old vertex \dots, then all vertices of this class are localized in $\square_1$''
\end{tabular}
\end{center}

The order is S11 (\cref{sup3:ssec:S11}), the properties of the cube split (\cref{sup3:ssec:S12}), S13 (\cref{sup3:ssec:S13}), the expansion S14 (\cref{sup3:ssec:S14a,sup3:ssec:S14b}) and S15 (\cref{sup3:ssec:S15}). Everything is at $L=5$, $d=3$, on the scale-$k$ lattice with $k\ge1$; the geometric lemmas hold for every $d$ with the same proofs. The step S10 (the regrouping after the separated lines are cut) is the main paper's; its cluster form is recalled below, because S11 and S15 are statements about the components it leaves. The chains without an old vertex (S16), the gauge covariance and the gauge-out of $B_0$ (S17--S18), the cube propagator (S19), the norms (S20), the renormalisation (S21) and the closure lemmas of the class family are in the fourth part.

\subsection*{Conventions for this part}

\paragraph{Scales and fields.} As in \cref{sup1:ssec:setting}: $\eta=L^{-k}$, $t=L^k\eps$, $e_k=e\,t^{1/2}$; $p(t)=b_0(1+\ln t^{-1})^{\bar p}$ with $\bar p>2$ and $r(t)=R(1+\ln t^{-1})^{\varrho}$ with $\varrho>1$, $R>R_0$ \cite[(2.3), (2.7)]{Ba2}; $q$ is the fixed real antisymmetric $N\times N$ matrix, $U(A_b)\defeq\exp(qe_k\eta A_b)$, and for an oriented contour $\Gamma$, $A(\Gamma)\defeq\sum_{b\subset\Gamma}\eta A_b$ and $U(A(\Gamma))\defeq\exp(qe_kA(\Gamma))$; all $U$ are orthogonal and commute, $U(A+A')=U(A)U(A')$. The covariant derivative is $(D^\eta_A\phi)(b)\defeq\eta^{-1}(U(A_b)\phi(b_+)-\phi(b_-))$ on the bond $b=(b_-,b_+)$, $Q_k(A)\phi(y)\defeq\sum_{x\in B^k(y)}\eta^dU(A(\Gamma^{(k)}_{y,x}))\phi(x)$ with the contours $\Gamma^{(k)}_{y,x}\subset B^k(y)$ of $\eta$-length $\le d$ \cite[(1.13)]{Ba3}, $P_k(A)=Q_k(A)^*Q_k(A)$, and on a union $X$ of unit blocks $H_X(A)\defeq-\Delta^{\eta,N}_{A,X}+m^2t^2+a_kP_k(A)$, $G_k(X,A)\defeq H_X(A)^{-1}$, with $-\Delta^{\eta,N}_{A,X}=(D^\eta_A)^*D^\eta_A$ over the bonds of $X$ and $m^2t^2>0$. The hypotheses of \cite[p.~412]{Ba3} on $\tilde B$ are those of \cref{sup1:ssec:setting}: $|\partial^\eta_\nu\tilde B_\mu|\le c_1p(t)$ on every pair of bonds of $\Omega$, and the H\"older quotients of exponent $\alpha_0$ of $\partial^\eta\tilde B$ are $\le c_1p(t)$ on pairs of points at distance $\le1$ at depth $\ge R_{\mathrm{Ba}}+2$ in $\Omega$, where $R_{\mathrm{Ba}}\ge R_0$ is the depth parameter of the background regularity of the main paper. The constant $c_1$ is independent of $k$ (this is the $k$-uniformity of the background bounds proved in the main paper's Step~3(C)); nothing else about $\tilde B$ is used.

\paragraph{Grid boxes and label boxes.} $|\cdot|_\infty$ is the sup norm; ``side'', ``diameter'' and $\dist$ are sup-norm notions unless said otherwise. A \emph{grid box} is a closed set $\prod_\mu[\alpha_\mu,\beta_\mu]\subset\R^d$ with all $\alpha_\mu\le\beta_\mu$ in the block-corner grid (a translate of $\Z^d$). A \emph{label box} is a grid box with all sides $\le2$. The localisation of \cite[p.~420]{Ba3} is read as in the main paper (its convention (Loc)): a label is attached to each leg of a vector field (the support box of its partition function $H_y\subset y+[-1,1]^d$, the partition $\{H_y\}_{y\in\text{grid}}$ being smooth with $\supp H_y\subset y+[-\frac34,\frac34]^d$ and derivatives bounded independently of $\eta$; a bond $b$ with $H_y(m_b)\ne0$ at its midpoint $m_b$ has both ends in $y+[-1,1]^d$), to each \cite[(1.6)]{Ba3} or \cite[(1.7)]{Ba3} vertex (the unit grid cube of its block), and to each block factor \cite[(1.12)--(1.15)]{Ba3} (the unit grid cube containing $B^k(y)$, which contains the contours $\Gamma^{(k)}_{y,x}$). The \emph{localisation box} $Q_v$ of a vertex $v$ is the intersection of the boxes of its labels; it is a label box (possibly degenerate) containing every summation point of the vertex. A counterterm vertex carries the labels of the anchor vertex of the subgraph it replaces (the main paper's reading (Agr)), so its localisation box contains its point. The distance ``$\dist(\square(v),\square(v'))$'' of \cite[pp.~432--433]{Ba3} is $\dist(Q_v,Q_{v'})$; the point-set distance of the summation points is at least this, so the rule ``distance $\ge1$ $\Rightarrow$ external field'' is applied only where the separation that \cite[(3.1)]{Ba3} needs holds for the actual points.

\paragraph{Classes and the cut rule.} The classes are the main paper's per-label classes: all graphs of a type and order with a given multiset $\mathcal B$ of label boxes, with the counterterm graphs of the divergent subgraphs; every member of a class has the same $\mathcal B$. The \emph{clusters} of $\mathcal B$ are the connected components of the graph on $\mathcal B$ in which two boxes are joined when they intersect, and $U_K$ is the union of the boxes of the cluster $K$. The cut rule of \cite[p.~432]{Ba3} in the main paper's form: a line is \emph{cut} (made an external field by \cite[(3.1)]{Ba3}, \cref{sup2:ssec:S9}) when its two end vertices lie in different clusters, and \emph{kept} otherwise. Every vertex lies in the intersection of its label boxes, so its labels lie in one cluster and ``the cluster of a vertex'' is well defined. A \emph{component} is a connected component of a graph with respect to its kept lines.

\paragraph{Vertex counts.} Three counts occur. The \emph{merged} count $V_{\mathrm{mer}}$ counts a product \cite[(1.18)]{Ba3} of two block factors at a common $y$ as one vertex; Ba{\l}aban's count $V_{\mathrm{Bal}}\le2V_{\mathrm{mer}}$ counts each of \cite[(1.13)--(1.15)]{Ba3} separately (\cite[p.~414]{Ba3}; this is the count in which degrees are computed, \cref{sup2:sec}); the \emph{label} count $V_{\mathrm{lab}}\le(2\bar n+1)V_{\mathrm{Bal}}$ is the number of label boxes, since a vertex has at most $2\bar n+1$ labels (at most $n+n'\le2\bar n+1$ vector legs for \cite[(1.8)--(1.11)]{Ba3}, the $R$-vertices having $n'=\bar n+1$ and $n\le\bar n$; one unit cube for \cite[(1.6)--(1.7)]{Ba3}; one unit cube and $n\le\bar n$ legs for a block factor). \Cref{sup3:lem:S113} bounds $V_{\mathrm{mer}}$ by a function of $\bar n$; with it $V_{\mathrm{lab}}(\bar n)\defeq(2\bar n+1)\cdot2\bar n(4N_A+5)$ and
\[
l_N\defeq2V_{\mathrm{lab}}(\bar n)
\]
are the constants of the main paper. Here $N_A$ is the fixed maximal number of insertions in the expansion of a scalar line or external scalar field in the fluctuation of the background, \cite[(3.45)]{Ba1} to order $\bar n$ on a line ($N_A=\bar n$) and the fixed odd order $n_{\mathrm{loop}}\ge\bar n$ of \cite[(3.47)]{Ba1} on a closed scalar loop; $N_A\defeq\max\{\bar n,n_{\mathrm{loop}}\}$ depends on $\bar n$ only.

\paragraph{The cube step.} $r_\square$ denotes the side parameter of the cube step: $r_\square=r^\flat$, the constant of the main paper's block-factor-$5$ reading (its cube step replaces the cubes of sides $r(t)$, $3r(t)$ by concentric cubes of constant sides $r^\flat$, $3r^\flat$), and $r_\square=r(t)$ wherever the printed step is applied as printed. $M$ is the side of the large blocks. Rounding a set outward to large blocks enlarges each side by at most $2M$.
